\documentclass[12pt]{article}
\usepackage[backend=biber,natbib=true,style=alphabetic,maxbibnames=50]{biblatex}
\addbibresource{/home/nqbh/reference/bib.bib}
\usepackage[utf8]{vietnam}
\usepackage{tocloft}
\renewcommand{\cftsecleader}{\cftdotfill{\cftdotsep}}
\usepackage[colorlinks=true,linkcolor=blue,urlcolor=red,citecolor=magenta]{hyperref}
\usepackage{amsmath,amssymb,amsthm,enumitem,fancyvrb,float,graphicx,mathtools,soul,tikz}
\usetikzlibrary{angles,calc,intersections,matrix,patterns,quotes,shadings}
\allowdisplaybreaks
\newtheorem{assumption}{Assumption}
\newtheorem{baitoan}{Bài toán}
\newtheorem{cauhoi}{Câu hỏi}
\newtheorem{conjecture}{Conjecture}
\newtheorem{corollary}{Corollary}
\newtheorem{dangtoan}{Dạng toán}
\newtheorem{definition}{Definition}
\newtheorem{dinhly}{Định lý}
\newtheorem{dinhnghia}{Định nghĩa}
\newtheorem{example}{Example}
\newtheorem{ghichu}{Ghi chú}
\newtheorem{hequa}{Hệ quả}
\newtheorem{hypothesis}{Hypothesis}
\newtheorem{lemma}{Lemma}
\newtheorem{luuy}{Lưu ý}
\newtheorem{nhanxet}{Nhận xét}
\newtheorem{notation}{Notation}
\newtheorem{note}{Note}
\newtheorem{principle}{Principle}
\newtheorem{problem}{Problem}
\newtheorem{proposition}{Proposition}
\newtheorem{question}{Question}
\newtheorem{remark}{Remark}
\newtheorem{theorem}{Theorem}
\newtheorem{vidu}{Ví dụ}
\usepackage[margin=2cm]{geometry}
\def\labelitemii{$\circ$}
\DeclareRobustCommand{\divby}{%
    \mathrel{\vbox{\baselineskip.65ex\lineskiplimit0pt\hbox{.}\hbox{.}\hbox{.}}}%
}
\setlist[itemize]{leftmargin=5mm}
\setlist[enumerate]{leftmargin=*}

\title{Fundamentals of Programming: Midterm Exam 2026\\Đề Thi Giữa Kỳ Cơ Sở Lập Trình 2026}
\author{Nguyễn Quản Bá Hồng}
\date{\today}

\begin{document}
\maketitle

\begin{center}
    \Large
    Tổng quan đề thi
\end{center}

\begin{table}[H]
    \small
    \centering
    \begin{tabular}{|c|l|c|c|c|}
        \hline
        {\bf STT} & {\bf Tên bài} & {\bf Thời gian} & {\bf Giới hạn bộ nhớ} & {\bf Điểm} \\
        \hline
        1 & Tổng hiệu tích thương & 1 giây & 1 GB & 100 \\
        \hline
        2 & Chu vi, diện tích hình chữ nhật & 1 giây & 1 GB & 100 \\
        \hline
        3 & Số nghiệm của phương trình bậc nhất & 1 giây & 1 GB & 100 \\
        \hline
        4 & Số nghiệm của phương trình bậc 2 & 1 giây & 1 GB & 100 \\
        \hline
        5 & Đếm chẵn lẻ & 1 giây & 1 GB & 100 \\
        \hline
    \end{tabular}
\end{table}
{\bf Lưu ý.} Thí sinh không cần phải pass hết bộ test mới được tính điểm, mà chỉ cần pass nhiều bộ test nhất có thể thì vẫn được tính điểm trên các bộ test đã được Online Judge chấp nhận là đúng.

\begin{baitoan}[Tổng hiệu tích thương]
    Tính tổng, hiệu, tích, thương của 2 số nguyên $a,b$.
    \item {\sf Input.} Mỗi test gồm 1 dòng duy nhất chứa 2 số nguyên $a,b$ ($-2\cdot10^9\le a,b\le2\cdot10^9$).
    \item {\sf Output.} Với mỗi test, in ra kết quả trên 1 dòng gồm 4 số nguyên, cách nhau bởi 1 khoảng trắng, lần lượt là tổng $a + b$, hiệu $a - b$, tích $ab$, và thương nguyên (integer division, sử dụng toán tử {\tt/} cho kiểu dữ liệu số nguyên, e.g., {\tt int, long long}, trong {\sf C++}), nếu $b = 0$ thì vẫn xuất ra tổng, hiệu, tích nhưng với thương thì xuất ra chuỗi {\tt division by zero}.
    \item {\sf Sample.}
    \begin{table}[H]
        \centering
        \begin{tabular}{|l|l|}
            \hline
            \verb|basic_arithmetic_operation.inp| & \verb|basic_arithmetic_operation.out| \\
            \hline
            1 0 & 1 1 0 division by zero \\
            \hline
            1 1 & 2 0 1 1\\
            \hline
            1 2 & 3 -1 2 0\\
            \hline
        \end{tabular}
    \end{table}
\end{baitoan}

\begin{baitoan}[Chu vi, diện tích hình chữ nhật]
    Tính chu vi, diện tích của 1 hình chữ nhật khi biết chiều dài \& chiều rộng.
    \item {\sf Input.} Mỗi test chứa gồm 1 dòng chứa 2 số nguyên dương $a,b$ ($-2\cdot10^9\le a,b\le2\cdot10^9$).
    \item {\sf Output.} Với mỗi test, in ra kết quả trên 1 dòng gồm 2 số nguyên, cách nhau bởi 1 khoảng trắng, lần lượt là chu vi \& diện tích của hình chữ nhật, nếu không tồn tại hình chữ nhật với chiều dài \& chiều rộng đã cho thì in ra {\tt-1}.
    \item {\sf Sample.}
    \begin{table}[H]
        \centering
        \begin{tabular}{|l|l|}
            \hline
            \verb|rectangle.inp| & \verb|rectangle.out| \\
            \hline
            1 1 & 4 1 \\
            \hline
            -1 1 &  -1\\
            \hline
        \end{tabular}
    \end{table}
\end{baitoan}

\begin{baitoan}[Số nghiệm của phương trình bậc nhất]
    Đếm số nghiệm của phương trình bậc nhất $ax + b = 0$ với $a, b$ là 2 số nguyên được nhập vào.
    \item {\sf Input.} Mỗi test gồm 1 dòng chứa duy nhất 2 số nguyên $a,b$ ($-2\cdot10^9\le a,b\le2\cdot10^9$).
    \item {\sf Output.} Với mỗi test, in ra kết quả trên 1 dòng. Nếu phương trình vô nghiệm thì in ra $0$, nếu phương trình vô số nghiệm thì in ra chuỗi {\tt infinity}, còn các trường hợp khác thì xuất ra số nghiệm của phương trình.
    \item {\sf Sample.}
    \begin{table}[H]
        \centering
        \begin{tabular}{|l|l|}
            \hline
            \verb|first_order_equation.inp| & \verb|first_order_equation.out| \\
            \hline
            1 1 & 1 \\
            \hline
            0 1 & 0 \\
            \hline
            0 0  & infinity \\
            \hline
        \end{tabular}
    \end{table}
\end{baitoan}

\begin{baitoan}[Số nghiệm của phương trình bậc 2]
    Đếm số nghiệm thực của phương trình bậc 2 $ax^2 + bx + c = 0$ với $a, b, c$ là 3 số nguyên được nhập vào.
    \item {\sf Input.} Mỗi test gồm 1 dòng chứa duy nhất 3 số nguyên $a,b,c$ ($-10^9\le a,b\le10^9$).
    \item {\sf Output.} Với mỗi test, in ra kết quả trên 1 dòng. Nếu phương trình vô nghiệm thì in ra $0$, nếu phương trình vô số nghiệm thì in ra chuỗi {\tt infinity}, còn các trường hợp khác thì xuất ra số nghiệm của phương trình.
    \item {\sf Sample.}
    \begin{table}[H]
        \centering
        \begin{tabular}{|l|l|}
            \hline
            \verb|second_order_equation.inp| & \verb|second_order_equation.out| \\
            \hline
            0 0 1 & 0 \\
            \hline
            0 1 1 & 1 \\
            \hline
            0 0 0 & infinity \\
            \hline
            1 2 10 & 0 \\
            \hline
            1 2 1 & 1 \\
            \hline
            1 3 2 & 2 \\
            \hline
        \end{tabular}
    \end{table}
\end{baitoan}

\begin{baitoan}[Đếm chẵn lẻ]
    Nhập vào $n$ số nguyên. Đếm số số chẵn, số số lẻ trong $n$ số nguyên đó.
    \item {\sf Input.} Mỗi test gồm 2 dòng. Dòng 1 chứa duy nhất 1 số nguyên dương $n$ ($1\le n\le2\cdot10^5$). Dòng 2 chứa $n$ số nguyên $a_1,a_2,\ldots,a_n$ ($-2\cdot10^9\le a_i\le2\cdot10^9$, $\forall i = 1,2,\ldots,n$).
    \item {\sf Output.} Với mỗi test, in ra số số chẵn \& số số lẻ trong $n$ số đó, cách nhau bởi 1 khoảng trắng.
    \item {\sf Sample.}
    \begin{table}[H]
        \centering
        \begin{tabular}{|l|l|}
            \hline
            \verb|count_even_odd.inp| & \verb|count_even_odd.out| \\
            \hline
            3 & 1 2 \\
            1 2 3 & \\
            \hline
            4 & 2 2 \\
            -1 2 -3 4 & \\
            \hline
        \end{tabular}
    \end{table}
\end{baitoan}

%------------------------------------------------------------------------------%

\printbibliography[heading=bibintoc]

\end{document}